\section{Experimental Setup}
\label{sec:experiments}

We designed experiments to test four sub-hypotheses addressing friction-completeness correlation, mechanism validation, gradient effects, and within-platform causal evidence.

\subsection{Sub-Hypothesis H1: Cross-Platform Correlation}

\textbf{Claim:} Higher friction score correlates with higher optional metadata field presence.

\textbf{Test:} Compare presence rates for 3 optional fields (preprocessing\_code, data\_source\_url, collection\_date) across platforms with friction scores 0 (UCI), 2 (OpenML), 3 (HuggingFace).

\textbf{Statistical Method:} Chi-squared tests for proportions, Cohen's $h$ effect sizes, Bonferroni correction ($\alpha=0.05/9=0.0056$) for 9 comparisons (3 fields $\times$ 3 platform pairs).

\textbf{Success Criteria:} Statistical significance ($p<0.0056$) AND practical significance (effect size $h>0.5$ medium effect) for $\geq$2/3 optional fields.

\subsection{Sub-Hypothesis H2: Enforcement Mechanism}

\textbf{Claim:} Required fields show high presence ($>$75\%) across ALL platforms, with weaker friction effects than optional fields.

\textbf{Test:} Compare presence rates for 2 required fields (license, version) across platforms, expecting high baseline ($>$75\%) even at friction=0.

\textbf{Statistical Method:} Chi-squared tests, Cramér's $V$ effect sizes, comparison with optional field effect sizes.

\textbf{Success Criteria:} All platforms show $>$75\% presence AND friction effect weaker than optional fields (smaller effect sizes).

\subsection{Sub-Hypothesis H3: Friction Gradient Effects}

\textbf{Claim:} Metadata presence increases monotonically with friction score (UCI $<$ OpenML $<$ HuggingFace) for fields where friction features apply.

\textbf{Test:} Test monotonic ordering for each field separately, expecting gradient effects for data\_source\_url and collection\_date (generic text fields benefiting from templates), potential threshold effects for preprocessing\_code (requires platform code snippet support).

\textbf{Statistical Method:} Pairwise chi-squared tests with directional hypotheses, Jonckheere-Terpstra trend test for monotonic ordering.

\textbf{Success Criteria:} Monotonic ordering confirmed for $\geq$2/3 optional fields, with documented explanation for deviations.

\subsection{Sub-Hypothesis H4: Within-Platform Causal Evidence}

\textbf{Claim:} Within HuggingFace platform, API-uploaded datasets show higher metadata completeness than manual-uploaded datasets.

\textbf{Test:} Compare mean completeness scores (0-100\%, calculated as mean presence across 6 fields) between 200 API-uploaded and 200 manual-uploaded datasets.

\textbf{Statistical Method:} Welch $t$-test (two-sided), normality checks via Shapiro-Wilk, Cohen's $d$ effect size, 95\% confidence intervals.

\textbf{Success Criteria:} Statistically significant difference ($p<0.05$) AND practical significance (Cohen's $d>0.5$ medium effect).

\subsection{Experimental Controls}

\textbf{Confound Control:}
\begin{itemize}
\item \textbf{Temporal drift:} Single-timepoint extraction (2026-08-19 snapshot)
\item \textbf{Platform size:} Stratified random sampling proportional to platform size
\item \textbf{Platform age/funding:} Within-platform design (H4) controls platform-level confounds
\item \textbf{Field selection:} Cover both optional (H1, H3) and required (H2) fields
\end{itemize}

\textbf{Statistical Power:} With $n=10{,}000$ total, $\alpha=0.05$, effect size $d=0.5$ (medium), power $>0.99$ for detecting 20pp+ differences. Within-platform test ($n=400$) provides power $>0.90$ for detecting $d=0.5$ effects.

\textbf{Multiple Comparisons:} Bonferroni correction applied to H1 (9 comparisons: 3 fields $\times$ 3 platform pairs), adjusted $\alpha=0.0056$. H2-H4 test distinct hypotheses with separate $\alpha=0.05$.
